\section*{Chapter \ref{ch:b2:cell-signalling} --- Chemical Messengers and Signal Transduction}

\begin{solution}{exo:b2:cell-signalling:1}
Insulin: endocrine, water-soluble (peptide). Cortisol: endocrine,
lipid-soluble. Acetylcholine at a synapse: synaptic, water-soluble.
Histamine at a wound: paracrine, water-soluble. Oestradiol:
endocrine, lipid-soluble. Adrenaline from the adrenal: endocrine,
water-soluble. Nitric oxide: paracrine, lipid-soluble (a gas that
crosses membranes).
\end{solution}

\begin{solution}{exo:b2:cell-signalling:2}
Adrenaline binds the receptor (off: dissociation, receptor
internalisation); receptor activates \omterm{prop:b2:cell-signalling:gpcr}{G protein}, GDP for GTP (off:
GTP hydrolysis in seconds); \omterm{prop:b2:cell-signalling:gpcr}{G protein} activates adenylyl cyclase,
which makes cAMP (off: phosphodiesterase); cAMP activates kinase A
(off: cAMP loss); kinase A phosphorylates phosphorylase kinase,
which phosphorylates glycogen phosphorylase (off: phosphatases);
phosphorylase releases glucose-1-phosphate from glycogen, converted
and exported as glucose.
\end{solution}

\begin{solution}{exo:b2:cell-signalling:3}
Surface receptors bind water-soluble messengers at the membrane and
act in seconds through \omterm{prop:b2:cell-signalling:gpcr}{second messengers} and kinases, changing the
activity of existing proteins. \omterm{prop:b2:cell-signalling:nuclear}{Nuclear receptors} bind lipid-soluble
messengers inside the cell and act in hours as \omterm{prop:b2:cell-differentiation:transcriptional}{transcription
factors}, changing which proteins are made.
\end{solution}

\begin{solution}{exo:b2:cell-signalling:4}
\omterm{prop:b2:cell-signalling:gpcr}{Cyclic AMP}: made by adenylyl cyclase, destroyed by
phosphodiesterase. Calcium: enters through channels or leaves the
endoplasmic reticulum in answer to inositol trisphosphate, removed
by pumps. Inositol trisphosphate (with diacylglycerol): made by
phospholipase C from a membrane lipid, removed by phosphatases.
\end{solution}

\begin{solution}{exo:b2:cell-signalling:5}
$\theta = [L]/(K_d + [L])$: $0.09$, $0.5$, $0.91$, $0.99$. Doubling the
receptors leaves the fraction unchanged and doubles the number
occupied --- and so the response.
\end{solution}

\begin{solution}{exo:b2:cell-signalling:6}
$50\times 200\times 100\times 300 = 3\times 10^{8}$. A hundred occupied
receptors give $3\times 10^{10}$ product molecules a second --- an
upper bound, since the stages saturate.
\end{solution}

\begin{solution}{exo:b2:cell-signalling:7}
The locked \omterm{prop:b2:cell-signalling:gpcr}{G protein} keeps the cyclase on, cAMP stays high, kinase A
keeps the chloride channel of the gut cells phosphorylated and
open, chloride is secreted, and sodium and water follow it into the
lumen --- litres a day. The lock is a covalent modification of the \omterm{prop:b2:cell-signalling:gpcr}{G
protein}, undone only when the cells themselves are replaced, days
later.
\end{solution}

\begin{solution}{exo:b2:cell-signalling:8}
$0.9\times 10^{-6}\,\text{mol/L}\times 2\times 10^{-12}\,\text{L} =
1.8\times 10^{-18}$ mol, $1.1\times 10^{6}$ ions (more in fact, since
buffers bind most of what enters). One channel at $10^{6}$ a second
would take a second; a thousand channels do it in a millisecond,
which is the time scale a signal needs.
\end{solution}

\begin{solution}{exo:b2:cell-signalling:9}
The two pathways converge on the same enzymes: glucagon's kinase A
phosphorylates glycogen synthase (off) and phosphorylase kinase
(on); insulin's cascade activates the phosphatase that removes the
phosphates. The enzymes' state is set by the balance of kinase and
phosphatase activity, so what counts is the ratio of the two
\omterm{def:b2:cell-signalling:kinds}{hormones}; with both high the phosphatase usually wins and glycogen
is stored.
\end{solution}

\begin{solution}{exo:b2:cell-signalling:10}
Adrenaline: cAMP diffuses across a \qty{20}{\micro m} cell in under a
second ($x^{2}/2D$ with $D \approx \qty{3e-10}{m^2/s}$), and every
enzyme of the cascade already exists --- seconds. Cortisol: a
transcript takes \qty{60}{s}, a protein \qty{100}{s}; from a few
dozen mRNAs each carrying a dozen ribosomes the cell makes a few
hundred enzyme molecules a minute, so the thousands needed for a
measurable effect take hours, after the minutes cortisol needs to
enter and reach the DNA.
\end{solution}

\begin{solution}{exo:b2:cell-signalling:11}
The receptor signals continuously: the MAP kinase cascade drives the
cell to divide without any growth factor, and it ignores the
absence of the signal that normally restrains it. Since a single
\omterm{def:b2:genome-diversification:mutation}{mutation} removes a control on division, such receptors (and the
kinases downstream) are among the commonest oncogenes. A drug that
occupies the kinase's ATP site stops the phosphorylations and
silences the receptor --- the principle of several cancer drugs.
\end{solution}

\begin{solution}{exo:b2:cell-signalling:12}
Endocrine: reaches every cell in a minute, cheap per cell, cannot
address one cell, acts for minutes to days --- indispensable for
whole-body states (fasting, growth, reproduction). Nervous:
millisecond, one target, expensive wiring, acts for milliseconds
--- indispensable for movement and rapid reflexes. Both are used for
the same message where speed and reach are both wanted: the
sympathetic nerves and the adrenal medulla both deliver
noradrenaline or adrenaline in a fright.
\end{solution}

\begin{solution}{pb:b2:cell-signalling:1}
\textbf{1.} $1\,\text{nmol}/5\,\text{L} = \qty{0.2}{nmol/L}$.
\textbf{2.} $\theta = 0.2/1.2 = 0.17$: about 3300 occupied receptors.
\textbf{3.} Three half-lives: \qty{0.025}{nmol/L}; $\theta = 0.024$.
\textbf{4.} About half a minute to a minute, the time the \omterm{def:b2:blood-circulation:blood}{blood}
takes to go round; the nerve delivers its transmitter directly to
the liver within a second.
\textbf{5.} The same fraction, but $33\,000$ occupied receptors: a
tenfold larger response.
\textbf{6.} At \qty{1}{nmol/L} a receptor with $K_d = \qty{1}{\micro
mol/L}$ is a thousandth occupied: no response. Matching $K_d$ to the
circulating range puts the steep part of the occupancy curve where
the \omterm{def:b2:cell-signalling:kinds}{hormone} actually varies.
\textbf{7.} $20\times 100\times 1\times 50\times 50\times 1000 =
10^{8}$ glucose molecules a second per occupied receptor.
\textbf{8.} $3300\times 20\times 100 = 6.6\times 10^{6}$ cAMP in the
first second; in \qty{5e-12}{L}, $1.1\times 10^{-17}$ mol, about
\qty{2}{\micro mol/L}.
\textbf{9.} $3300\times 10^{8} = 3.3\times 10^{11}$ a second, $2\times
10^{13}$ a minute per cell.
\textbf{10.} $2\times 10^{11}$ cells $\times 2\times 10^{13} = 4\times
10^{24}$ molecules a minute, 6.6 mol, \qty{1200}{g} a minute --- two
hundred times the observed \qty{5}{g/min}. The gains multiply only
while every stage is unsaturated; in fact the phosphorylase is
limited by its substrate and by the phosphatases, and the export of
glucose by transporters. The cascade's gain is spent on speed and
sensitivity, not on output.
\textbf{11.} At \qty{5}{g/min}, twenty minutes; the fright is over in
minutes and the off-switches stop the cascade long before.
\textbf{12.} GTP hydrolysis by the \omterm{prop:b2:cell-signalling:gpcr}{G protein} (seconds);
phosphodiesterase on cAMP (seconds); phosphatases on the
phosphorylated enzymes (seconds to minutes); \omterm{prop:b2:reproductive-hormones:pulses}{receptor
desensitisation} (minutes). The fastest are the \omterm{prop:b2:cell-signalling:gpcr}{G protein}'s own clock
and the phosphodiesterase.
\textbf{13.} cAMP is not destroyed: it accumulates and persists, the
kinase stays on, and the glucose output is larger and lasts longer
--- the alertness of coffee.
\textbf{14.} $4000/50 = \qty{80}{s}$.
\textbf{15.} $600/5 = \qty{120}{s}$ per enzyme per ribosome; 30 per
ribosome per hour, 600 per mRNA per hour.
\textbf{16.} $100\times 600 = 60\,000$ enzymes an hour; $10^{5}$ after
about \qty{1.7}{h}.
\textbf{17.} Ten to twenty minutes for cortisol to enter, bind and
reach its genes, an hour before the first mRNAs are numerous and
translated, then two hours of accumulation: about four hours.
\textbf{18.} A cascade changes the activity of enzymes that exist;
cortisol's job is to change which enzymes exist, for days, which
only transcription can do. A fright cannot wait four hours for new
enzymes.
\textbf{19.} More receptors means more occupied receptors at any
adrenaline concentration: the occupancy fraction is unchanged but
the response at each concentration is larger --- the curve is
scaled up, not shifted. A slow \omterm{def:b2:cell-signalling:kinds}{hormone} sets the gain of a fast one.
\textbf{20.} Glucose enters the $\beta$ cell and is metabolised; ATP
rises and closes an ATP-sensitive potassium channel; the membrane
depolarises; voltage-gated calcium channels open; calcium triggers
the exocytosis of insulin granules --- metabolism coupled to
secretion by an electrical step.
\textbf{21.} The two receptors start two pathways that converge on
the same enzymes: glucagon's kinase A phosphorylates them, insulin's
cascade activates the phosphatase that dephosphorylates them; the
enzymes' state, and so the direction of glycogen metabolism, follows
the ratio of the two.
\textbf{22.} Fasting glucose high (the liver, unrestrained, keeps
making it), post-meal glucose very high (no uptake by muscle and
fat); without insulin the fat stores release fatty acids, which the
liver converts to ketone acids --- ketoacidosis.
\textbf{23.} Working muscle takes up glucose without insulin, so on a
long run the glucose falls while the injected insulin, unregulated,
goes on suppressing the liver: hypoglycaemia and collapse.
\textbf{24.} Glucose loop: sensor the $\beta$ and $\alpha$ cells,
two effector \omterm{def:b2:cell-signalling:kinds}{hormones} for the two directions, a high gain holding
the level within a fifth, minutes to act. \omterm{def:b2:blood-pressure:baroreflex}{Baroreflex}: stretch
receptors, the \omterm{def:b2:heart:anatomy}{heart} and vessels as effectors, gain about four, a
second to act.
\textbf{25.} Occupancy $0.17$ after release; gain $10^{8}$ per second
per occupied receptor; liver output \qty{5}{g/min} in fact; cortisol
takes about four hours.
\end{solution}
